\section{Discussion and Mechanistic Insights}
\label{sec:discussion}

This section investigates the internal optimization dynamics of dual attention under SGD, examines data augmentation chirality pitfalls, evaluates deployment profiles, and discusses future directions.

\subsection{Optimization Diagnostics: The $\alpha/\gamma$ Gradient Deadlock}
\label{sec:deadlock}
A common convention in dual attention~\cite{fu2019dual,sun2024datransunet} is the zero-initialized residual scaling:
\begin{equation}
    \mathbf{Y}_{\text{PAM}} = \alpha \mathbf{E}_{\text{PAM}} + \mathbf{X}, \quad \alpha_{\text{init}} = 0,
\end{equation}
and similarly $\beta_{\text{init}} = 0$ for CAM. However, auditing gradient dynamics under standard SGD reveals a hidden failure mode: \textbf{the $\alpha/\gamma$ gradient deadlock}.

By the chain rule, the gradient with respect to projection weights $\mathbf{W}_Q$ is:
\begin{equation}
    \frac{\partial \mathcal{L}}{\partial \mathbf{W}_Q} = \frac{\partial \mathcal{L}}{\partial \mathbf{Y}_{\text{PAM}}} \cdot \left( \alpha \frac{\partial \mathbf{E}_{\text{PAM}}}{\partial \mathbf{W}_Q} \right).
\end{equation}
At initialization, because $\alpha = 0$, $\partial \mathcal{L} / \partial \mathbf{W}_Q$ evaluates strictly to zero. Moreover, $|\partial \mathcal{L} / \partial \alpha| = |\langle \partial \mathcal{L}/\partial \mathbf{Y}_{\text{PAM}}, \mathbf{E}_{\text{PAM}} \rangle|$ measures only $4 \times 10^{-6}$ to $4 \times 10^{-5}$, four orders of magnitude weaker than decoder convolutional gradients ($\approx 0.29$).

\paragraph{Empirical Verification} Tracking the ratio of maximum weight magnitude $\max |w|$ after 300 epochs relative to initialization bound $b$ shows that under pure weight decay with zero backpropagation, parameters scale down by $(1 - \eta \lambda)^T \approx 0.84$--$0.86$. In reproduced DA-TransUNet models, PAM projection weights exhibited drift ratios of exactly $0.84$--$0.86$, confirming that attention weights were never actively updated by gradient descent. Concurrently, forcing $\alpha = 0$ at test time altered fewer than $0.0002\%$ of predicted pixels.

\paragraph{The True Driver of Dual-Attention Blocks} This diagnostic demonstrates that the empirical success of dual attention in CNN-Transformer hybrids does not originate from full-rank spatial self-attention matrices. Instead, performance gains are driven primarily by the internal $3 \times 3$ bottleneck convolutional compression ($C \to C/16 \to C$). This insight explains why our tri-axial approximation ($M=28, r=32, G=8$) achieves full segmentation accuracy while shedding $90\%$ of model parameters.

\subsection{Anatomical Symmetries and Data Augmentation Pitfalls}
\label{sec:symmetry}
Decomposing segmentation errors on Synapse CT reveals that while $87\%$ of errors occur within $\le 2$ mm of true boundaries, $9\%$ consist of large block misidentifications ($>20$ px), primarily left/right kidney swaps. 

This error pattern stems directly from standard data augmentation: \texttt{random\_rot\_flip} applies random $90^\circ$ rotations followed by mandatory axis reflections without transposing bilateral organ labels. Consequently, half of all training slices exhibit reversed anatomical chirality (artificial situs inversus). Disabling mirror reflection eliminated left/right kidney confusion ($0.0\%$ error), though small training sets benefit from reflection to prevent shape overfitting. Clinicians and researchers should carefully deploy chirality-aware augmentations for bilateral anatomy.

\subsection{Computational Efficiency and Clinical Deployability}

\begin{table}[t]
\centering
\caption{Deployment profile comparison on a single NVIDIA RTX 4090 GPU (batch size 1, $224 \times 224$ input).}
\label{tab:deployment}
\resizebox{\columnwidth}{!}{
\begin{tabular}{lcc}
\toprule
\textbf{Metric} & \textbf{DA-TransUNet} & \textbf{ApproxDA-TransUNet (Ours)} \\
\midrule
Model Parameters & 107.95 M & \textbf{10.75 M} (\textbf{--90.0\%}) \\
Inference GFLOPs & 30.25 G & \textbf{6.92 G} (\textbf{--77.1\%}) \\
Peak GPU Memory (Training, $B{=}24$) & 16.8 GB & \textbf{6.1 GB} (\textbf{--63.7\%}) \\
Inference Latency (Single Slice) & 28.4 ms & \textbf{8.7 ms} (\textbf{3.3$\times$ Speedup}) \\
Throughput (Slices/sec) & 35.2 & \textbf{114.9} (\textbf{+226\%}) \\
Storage Footprint (FP32 Checkpoint) & 432 MB & \textbf{43 MB} \\
\bottomrule
\end{tabular}
}
\end{table}

As summarized in Table~\ref{tab:deployment}, ApproxDA-TransUNet delivers major hardware benefits:
\begin{itemize}
    \item \textbf{Memory \& Throughput}: Peak training memory drops by $63.7\%$ ($16.8 \to 6.1$ GB), enabling training on consumer GPUs. Inference latency drops from $28.4$ ms to $8.7$ ms ($114.9$ slices/s), achieving a $3.3\times$ speedup suitable for real-time intraoperative guidance.
    \item \textbf{Edge Deployment}: The storage footprint is reduced from $432$ MB to $43$ MB, facilitating edge deployment on embedded scanners.
\end{itemize}

\subsection{Limitations and Future Work}
ApproxDA-TransUNet processes 3D scans on a 2D planar basis. While highly efficient, approximately $12\%$ of remaining errors occur at the cranial/caudal extremities of organs where 3D inter-slice context is needed to identify organ endpoints. Extending ApproxDA blocks to lightweight 2.5D or anisotropic 3D formulations presents a promising avenue for volumetric research.
